\documentclass[11pt,a4paper]{article}
% ==============================================================
%  Préambule commun aux fiches -- Advanced Probabilities (M1)
%  Usage, dans chaque fiche :
%     \documentclass[11pt,a4paper]{article}
%     % ==============================================================
%  Préambule commun aux fiches -- Advanced Probabilities (M1)
%  Usage, dans chaque fiche :
%     \documentclass[11pt,a4paper]{article}
%     % ==============================================================
%  Préambule commun aux fiches -- Advanced Probabilities (M1)
%  Usage, dans chaque fiche :
%     \documentclass[11pt,a4paper]{article}
%     \input{../preambule.tex}
%     \fiche{C5}{Vecteurs gaussiens}{Gaussian vectors}
% ==============================================================
\usepackage[T1]{fontenc}
\usepackage[utf8]{inputenc}
\usepackage[french]{babel}
\usepackage{lmodern}
\usepackage[a4paper,margin=1.9cm,top=1.7cm,bottom=1.9cm]{geometry}
\usepackage{amsmath,amssymb,amsthm,mathtools}
\usepackage[inline]{enumitem}
\usepackage{booktabs,array,tabularx}
\usepackage{xcolor}
\usepackage[most]{tcolorbox}
\usepackage{fancyhdr}
\usepackage{hyperref}
\hypersetup{colorlinks=true,linkcolor=blue!50!black,urlcolor=blue!50!black}

\setlist{itemsep=1pt,topsep=3pt,parsep=0pt}
\setlength{\parindent}{0pt}
\setlength{\parskip}{3pt}
\emergencystretch=3em

% ---------- Identification de la fiche -------------------------
\newcommand{\ficheid}{}
\newcommand{\fichetitre}{}
\newcommand{\fiche}[3]{%
  \renewcommand{\ficheid}{#1}\renewcommand{\fichetitre}{#2}%
  \pagestyle{fancy}\fancyhf{}%
  \fancyhead[L]{\small\textsf{Advanced Probabilities -- Fiche #1}}%
  \fancyhead[R]{\small\textsf{#2}}%
  \fancyfoot[C]{\small\thepage}%
  \begin{center}
    {\LARGE\bfseries Fiche #1 -- #2}\par\smallskip
    {\large\itshape (#3)}\par\smallskip
    {\small Cours de V.\,M.~Lambert, M1 -- semestre 1}
  \end{center}\vspace{-2pt}\hrule\vspace{6pt}}

% ---------- Environnements numérotés (un seul compteur) --------
\theoremstyle{definition}
\newtheorem{definition}{Définition}
\renewcommand{\thedefinition}{\ficheid.\arabic{definition}}
\newtheorem{theoreme}[definition]{Théorème}
\newtheorem{proposition}[definition]{Proposition}
\newtheorem{lemme}[definition]{Lemme}
\newtheorem{corollaire}[definition]{Corollaire}
\newtheorem{exemple}[definition]{Exemple}
\theoremstyle{remark}
\newtheorem{remarque}[definition]{Remarque}

% Résultat hors cours (standard, signalé comme tel)
\theoremstyle{definition}
\newtheorem{horscours}[definition]{Hors cours}

% Idée de preuve
\newenvironment{idee}{\par\small\textit{Idée de preuve.}~}{\par}

% ---------- Encadrés -----------------------------------------
\tcbset{fichebox/.style={enhanced,breakable,boxrule=0.6pt,arc=2pt,
  left=5pt,right=5pt,top=3pt,bottom=3pt,fonttitle=\bfseries}}
\newtcolorbox{piege}[1][]{fichebox,colback=red!4,colframe=red!60!black,
  title={Pièges classiques},#1}
\newtcolorbox{liens}[1][]{fichebox,colback=blue!3,colframe=blue!50!black,
  title={Liens entre notions},#1}
\newtcolorbox{methode}[1][Méthode]{fichebox,colback=green!4,colframe=green!40!black,title={#1}}
\newtcolorbox{alerte}[1][]{fichebox,colback=orange!6,colframe=orange!70!black,
  title={Attention à l'énoncé},#1}

% Exercice de TD (fiches de rappels)
\newcounter{exo}
\newenvironment{exo}[1][]{\refstepcounter{exo}\par\medskip
  \noindent{\large\bfseries Exercice~\theexo}\ifx\relax#1\relax\else~{\large\bfseries -- #1}\fi\par\nopagebreak\smallskip}{\par}
\newcommand{\outils}{\textbf{Outils du cours.}~}
\newcommand{\meth}{\textbf{Méthode / astuce.}~}
\newcommand{\indic}{\textbf{Indication.}~}
\newcommand{\rep}{\textbf{Réponse.}~}

% ---------- Notations du cours ---------------------------------
\newcommand{\R}{\mathbb{R}}
\newcommand{\N}{\mathbb{N}}
\newcommand{\Q}{\mathbb{Q}}
\newcommand{\C}{\mathbb{C}}
\newcommand{\PP}{\mathbb{P}}
\newcommand{\QQ}{\mathbb{Q}}
\newcommand{\E}{\mathbb{E}}
\newcommand{\F}{\mathcal{F}}
\newcommand{\G}{\mathcal{G}}
\newcommand{\Hc}{\mathcal{H}}
\newcommand{\Tc}{\mathcal{T}}
\newcommand{\Bc}{\mathcal{B}}
\newcommand{\Cc}{\mathcal{C}}
\newcommand{\Ec}{\mathcal{E}}
\newcommand{\Pc}{\mathcal{P}}
\newcommand{\Nc}{\mathcal{N}}
\newcommand{\LN}{\mathcal{LN}}
\newcommand{\Expo}{\mathcal{E}}
\newcommand{\Poi}{\mathcal{P}}
\newcommand{\Ff}{\mathbb{F}}
\newcommand{\1}{\mathbf{1}}
\newcommand{\dd}{\,\mathrm{d}}
\DeclareMathOperator{\Var}{\mathbb{V}}
\DeclareMathOperator{\Cov}{Cov}
\DeclareMathOperator{\card}{card}
\DeclareMathOperator{\diag}{diag}
\newcommand{\indep}{\perp\!\!\!\perp}
\newcommand{\ps}{\ \PP\text{-p.s.}}
\newcommand{\tendloi}{\xrightarrow[n\to+\infty]{\text{loi}}}
\newcommand{\tendps}{\xrightarrow[n\to+\infty]{\text{p.s.}}}
\newcommand{\tendproba}{\xrightarrow[n\to+\infty]{\PP}}
\newcommand{\tendLp}{\xrightarrow[n\to+\infty]{L^p}}
\newcommand{\tend}{\xrightarrow[n\to+\infty]{}}

%     \fiche{C5}{Vecteurs gaussiens}{Gaussian vectors}
% ==============================================================
\usepackage[T1]{fontenc}
\usepackage[utf8]{inputenc}
\usepackage[french]{babel}
\usepackage{lmodern}
\usepackage[a4paper,margin=1.9cm,top=1.7cm,bottom=1.9cm]{geometry}
\usepackage{amsmath,amssymb,amsthm,mathtools}
\usepackage[inline]{enumitem}
\usepackage{booktabs,array,tabularx}
\usepackage{xcolor}
\usepackage[most]{tcolorbox}
\usepackage{fancyhdr}
\usepackage{hyperref}
\hypersetup{colorlinks=true,linkcolor=blue!50!black,urlcolor=blue!50!black}

\setlist{itemsep=1pt,topsep=3pt,parsep=0pt}
\setlength{\parindent}{0pt}
\setlength{\parskip}{3pt}
\emergencystretch=3em

% ---------- Identification de la fiche -------------------------
\newcommand{\ficheid}{}
\newcommand{\fichetitre}{}
\newcommand{\fiche}[3]{%
  \renewcommand{\ficheid}{#1}\renewcommand{\fichetitre}{#2}%
  \pagestyle{fancy}\fancyhf{}%
  \fancyhead[L]{\small\textsf{Advanced Probabilities -- Fiche #1}}%
  \fancyhead[R]{\small\textsf{#2}}%
  \fancyfoot[C]{\small\thepage}%
  \begin{center}
    {\LARGE\bfseries Fiche #1 -- #2}\par\smallskip
    {\large\itshape (#3)}\par\smallskip
    {\small Cours de V.\,M.~Lambert, M1 -- semestre 1}
  \end{center}\vspace{-2pt}\hrule\vspace{6pt}}

% ---------- Environnements numérotés (un seul compteur) --------
\theoremstyle{definition}
\newtheorem{definition}{Définition}
\renewcommand{\thedefinition}{\ficheid.\arabic{definition}}
\newtheorem{theoreme}[definition]{Théorème}
\newtheorem{proposition}[definition]{Proposition}
\newtheorem{lemme}[definition]{Lemme}
\newtheorem{corollaire}[definition]{Corollaire}
\newtheorem{exemple}[definition]{Exemple}
\theoremstyle{remark}
\newtheorem{remarque}[definition]{Remarque}

% Résultat hors cours (standard, signalé comme tel)
\theoremstyle{definition}
\newtheorem{horscours}[definition]{Hors cours}

% Idée de preuve
\newenvironment{idee}{\par\small\textit{Idée de preuve.}~}{\par}

% ---------- Encadrés -----------------------------------------
\tcbset{fichebox/.style={enhanced,breakable,boxrule=0.6pt,arc=2pt,
  left=5pt,right=5pt,top=3pt,bottom=3pt,fonttitle=\bfseries}}
\newtcolorbox{piege}[1][]{fichebox,colback=red!4,colframe=red!60!black,
  title={Pièges classiques},#1}
\newtcolorbox{liens}[1][]{fichebox,colback=blue!3,colframe=blue!50!black,
  title={Liens entre notions},#1}
\newtcolorbox{methode}[1][Méthode]{fichebox,colback=green!4,colframe=green!40!black,title={#1}}
\newtcolorbox{alerte}[1][]{fichebox,colback=orange!6,colframe=orange!70!black,
  title={Attention à l'énoncé},#1}

% Exercice de TD (fiches de rappels)
\newcounter{exo}
\newenvironment{exo}[1][]{\refstepcounter{exo}\par\medskip
  \noindent{\large\bfseries Exercice~\theexo}\ifx\relax#1\relax\else~{\large\bfseries -- #1}\fi\par\nopagebreak\smallskip}{\par}
\newcommand{\outils}{\textbf{Outils du cours.}~}
\newcommand{\meth}{\textbf{Méthode / astuce.}~}
\newcommand{\indic}{\textbf{Indication.}~}
\newcommand{\rep}{\textbf{Réponse.}~}

% ---------- Notations du cours ---------------------------------
\newcommand{\R}{\mathbb{R}}
\newcommand{\N}{\mathbb{N}}
\newcommand{\Q}{\mathbb{Q}}
\newcommand{\C}{\mathbb{C}}
\newcommand{\PP}{\mathbb{P}}
\newcommand{\QQ}{\mathbb{Q}}
\newcommand{\E}{\mathbb{E}}
\newcommand{\F}{\mathcal{F}}
\newcommand{\G}{\mathcal{G}}
\newcommand{\Hc}{\mathcal{H}}
\newcommand{\Tc}{\mathcal{T}}
\newcommand{\Bc}{\mathcal{B}}
\newcommand{\Cc}{\mathcal{C}}
\newcommand{\Ec}{\mathcal{E}}
\newcommand{\Pc}{\mathcal{P}}
\newcommand{\Nc}{\mathcal{N}}
\newcommand{\LN}{\mathcal{LN}}
\newcommand{\Expo}{\mathcal{E}}
\newcommand{\Poi}{\mathcal{P}}
\newcommand{\Ff}{\mathbb{F}}
\newcommand{\1}{\mathbf{1}}
\newcommand{\dd}{\,\mathrm{d}}
\DeclareMathOperator{\Var}{\mathbb{V}}
\DeclareMathOperator{\Cov}{Cov}
\DeclareMathOperator{\card}{card}
\DeclareMathOperator{\diag}{diag}
\newcommand{\indep}{\perp\!\!\!\perp}
\newcommand{\ps}{\ \PP\text{-p.s.}}
\newcommand{\tendloi}{\xrightarrow[n\to+\infty]{\text{loi}}}
\newcommand{\tendps}{\xrightarrow[n\to+\infty]{\text{p.s.}}}
\newcommand{\tendproba}{\xrightarrow[n\to+\infty]{\PP}}
\newcommand{\tendLp}{\xrightarrow[n\to+\infty]{L^p}}
\newcommand{\tend}{\xrightarrow[n\to+\infty]{}}

%     \fiche{C5}{Vecteurs gaussiens}{Gaussian vectors}
% ==============================================================
\usepackage[T1]{fontenc}
\usepackage[utf8]{inputenc}
\usepackage[french]{babel}
\usepackage{lmodern}
\usepackage[a4paper,margin=1.9cm,top=1.7cm,bottom=1.9cm]{geometry}
\usepackage{amsmath,amssymb,amsthm,mathtools}
\usepackage[inline]{enumitem}
\usepackage{booktabs,array,tabularx}
\usepackage{xcolor}
\usepackage[most]{tcolorbox}
\usepackage{fancyhdr}
\usepackage{hyperref}
\hypersetup{colorlinks=true,linkcolor=blue!50!black,urlcolor=blue!50!black}

\setlist{itemsep=1pt,topsep=3pt,parsep=0pt}
\setlength{\parindent}{0pt}
\setlength{\parskip}{3pt}
\emergencystretch=3em

% ---------- Identification de la fiche -------------------------
\newcommand{\ficheid}{}
\newcommand{\fichetitre}{}
\newcommand{\fiche}[3]{%
  \renewcommand{\ficheid}{#1}\renewcommand{\fichetitre}{#2}%
  \pagestyle{fancy}\fancyhf{}%
  \fancyhead[L]{\small\textsf{Advanced Probabilities -- Fiche #1}}%
  \fancyhead[R]{\small\textsf{#2}}%
  \fancyfoot[C]{\small\thepage}%
  \begin{center}
    {\LARGE\bfseries Fiche #1 -- #2}\par\smallskip
    {\large\itshape (#3)}\par\smallskip
    {\small Cours de V.\,M.~Lambert, M1 -- semestre 1}
  \end{center}\vspace{-2pt}\hrule\vspace{6pt}}

% ---------- Environnements numérotés (un seul compteur) --------
\theoremstyle{definition}
\newtheorem{definition}{Définition}
\renewcommand{\thedefinition}{\ficheid.\arabic{definition}}
\newtheorem{theoreme}[definition]{Théorème}
\newtheorem{proposition}[definition]{Proposition}
\newtheorem{lemme}[definition]{Lemme}
\newtheorem{corollaire}[definition]{Corollaire}
\newtheorem{exemple}[definition]{Exemple}
\theoremstyle{remark}
\newtheorem{remarque}[definition]{Remarque}

% Résultat hors cours (standard, signalé comme tel)
\theoremstyle{definition}
\newtheorem{horscours}[definition]{Hors cours}

% Idée de preuve
\newenvironment{idee}{\par\small\textit{Idée de preuve.}~}{\par}

% ---------- Encadrés -----------------------------------------
\tcbset{fichebox/.style={enhanced,breakable,boxrule=0.6pt,arc=2pt,
  left=5pt,right=5pt,top=3pt,bottom=3pt,fonttitle=\bfseries}}
\newtcolorbox{piege}[1][]{fichebox,colback=red!4,colframe=red!60!black,
  title={Pièges classiques},#1}
\newtcolorbox{liens}[1][]{fichebox,colback=blue!3,colframe=blue!50!black,
  title={Liens entre notions},#1}
\newtcolorbox{methode}[1][Méthode]{fichebox,colback=green!4,colframe=green!40!black,title={#1}}
\newtcolorbox{alerte}[1][]{fichebox,colback=orange!6,colframe=orange!70!black,
  title={Attention à l'énoncé},#1}

% Exercice de TD (fiches de rappels)
\newcounter{exo}
\newenvironment{exo}[1][]{\refstepcounter{exo}\par\medskip
  \noindent{\large\bfseries Exercice~\theexo}\ifx\relax#1\relax\else~{\large\bfseries -- #1}\fi\par\nopagebreak\smallskip}{\par}
\newcommand{\outils}{\textbf{Outils du cours.}~}
\newcommand{\meth}{\textbf{Méthode / astuce.}~}
\newcommand{\indic}{\textbf{Indication.}~}
\newcommand{\rep}{\textbf{Réponse.}~}

% ---------- Notations du cours ---------------------------------
\newcommand{\R}{\mathbb{R}}
\newcommand{\N}{\mathbb{N}}
\newcommand{\Q}{\mathbb{Q}}
\newcommand{\C}{\mathbb{C}}
\newcommand{\PP}{\mathbb{P}}
\newcommand{\QQ}{\mathbb{Q}}
\newcommand{\E}{\mathbb{E}}
\newcommand{\F}{\mathcal{F}}
\newcommand{\G}{\mathcal{G}}
\newcommand{\Hc}{\mathcal{H}}
\newcommand{\Tc}{\mathcal{T}}
\newcommand{\Bc}{\mathcal{B}}
\newcommand{\Cc}{\mathcal{C}}
\newcommand{\Ec}{\mathcal{E}}
\newcommand{\Pc}{\mathcal{P}}
\newcommand{\Nc}{\mathcal{N}}
\newcommand{\LN}{\mathcal{LN}}
\newcommand{\Expo}{\mathcal{E}}
\newcommand{\Poi}{\mathcal{P}}
\newcommand{\Ff}{\mathbb{F}}
\newcommand{\1}{\mathbf{1}}
\newcommand{\dd}{\,\mathrm{d}}
\DeclareMathOperator{\Var}{\mathbb{V}}
\DeclareMathOperator{\Cov}{Cov}
\DeclareMathOperator{\card}{card}
\DeclareMathOperator{\diag}{diag}
\newcommand{\indep}{\perp\!\!\!\perp}
\newcommand{\ps}{\ \PP\text{-p.s.}}
\newcommand{\tendloi}{\xrightarrow[n\to+\infty]{\text{loi}}}
\newcommand{\tendps}{\xrightarrow[n\to+\infty]{\text{p.s.}}}
\newcommand{\tendproba}{\xrightarrow[n\to+\infty]{\PP}}
\newcommand{\tendLp}{\xrightarrow[n\to+\infty]{L^p}}
\newcommand{\tend}{\xrightarrow[n\to+\infty]{}}

\begin{document}
\fiche{C5}{Vecteurs gaussiens}{Gaussian vectors}

\textit{Source : slides, chapitre 6 (Déf.~35, 41, Lemmes 8--10, Thm~8--11) et poly « Gaussian vector » (Déf.~1, 5, Lemmes 6--8, 11,
Thm~9, 10, 12, 13). Convergence en loi, LGN, TCL : fiche C4.}

\section*{1. Gaussienne réelle}

\begin{definition}[Loi gaussienne sur $\R$]\label{C5:gauss1}
Une probabilité $\mu$ sur $\Bc(\R)$ est \emph{gaussienne} si
\begin{enumerate}
  \item (cas non dégénéré) $\mu=f_{m,\sigma^2}\cdot\lambda$ avec
  $f_{m,\sigma^2}(y)=\dfrac{1}{\sqrt{2\pi\sigma^2}}\,e^{-\frac{(y-m)^2}{2\sigma^2}}$, $m\in\R$, $\sigma^2\in\,]0,\infty[$ ;
  \item (cas dégénéré) $\mu=\delta_m$, masse de Dirac en $m$ : $\delta_m(A)=1$ si $m\in A$, $0$ sinon (« $\sigma^2=0$ »).
\end{enumerate}
Si $\PP_X=\mu$ : $\E(X)=\int x\dd\mu=m$, $\Var(X)=\sigma^2$, et $\varphi_X(t)=e^{itm-\frac{t^2\sigma^2}{2}}$. On note $X\sim\Nc(m,\sigma^2)$.
\end{definition}

\begin{proposition}[combinaisons de gaussiennes indépendantes]\label{C5:somme-indep}
Si $X_1,\dots,X_n$ sont \textbf{indépendantes}, $X_i\sim\Nc(m_i,\sigma_i^2)$, et $a_i,b_i\in\R$, alors
\[\sum_{i=1}^n(a_iX_i+b_i)\sim\Nc\Bigl(\sum_{i=1}^n(a_im_i+b_i),\ \sum_{i=1}^n(a_i\sigma_i)^2\Bigr).\]
En particulier ($n=1$, standard) $aX+b\sim\Nc(am+b,a^2\sigma^2)$ et $\frac{X-m}{\sigma}\sim\Nc(0,1)$ si $\sigma>0$.
\end{proposition}
\begin{idee}
Indépendance (C3) : $\varphi_{\sum(a_jX_j+b_j)}(t)=\prod_je^{itb_j}\varphi_{X_j}(a_jt)
=\exp\bigl(it\sum_j(a_jm_j+b_j)-\frac{t^2}2\sum_j(a_j\sigma_j)^2\bigr)$, et $\varphi$ caractérise la loi.
\end{idee}

\section*{2. Vecteurs gaussiens}

\begin{definition}[Vecteur gaussien]\label{C5:vecteur}
$X=(X_1,\dots,X_d)$ à valeurs dans $\R^d$ est un \emph{vecteur gaussien} si \textbf{toute} combinaison linéaire de ses
composantes est une v.a.\ réelle gaussienne : $\forall a\in\R^d$, $a^tX=\sum_{i=1}^da_iX_i$ est gaussienne (éventuellement dégénérée).
Une loi gaussienne sur $\R^d$ est la loi d'un vecteur gaussien.
\end{definition}

\begin{lemme}\label{C5:indep-vect}
Si $X_1,\dots,X_d$ sont des gaussiennes réelles \textbf{indépendantes}, $(X_1,\dots,X_d)$ est un vecteur gaussien, et
$\sum a_iX_i\sim\Nc\bigl(\sum a_im_i,\sum a_i^2\sigma_i^2\bigr)$.
\end{lemme}

\begin{lemme}[transformation affine]\label{C5:affine}
Soit $X$ un vecteur gaussien de $\R^d$, de moyenne $m=\E X=(\E X_1,\dots,\E X_d)^t$ et de matrice de variance-covariance
$\Sigma_X=(\Cov(X_i,X_j))_{i,j}=\E\bigl((X-\E X)(X-\E X)^t\bigr)$. Pour $A\in\mathcal M_{n,d}(\R)$, $b\in\R^n$,
$Y=AX+b$ est un vecteur gaussien de $\R^n$, de moyenne $Am+b$ et de covariance $A\Sigma_XA^t$.
\end{lemme}
\begin{idee}
Pour $a\in\R^n$, $a^tY=a^tb+(A^ta)^tX$ est gaussienne car $X$ est un vecteur gaussien.
\end{idee}

\begin{remarque}[standard]
$\Sigma_X$ est symétrique positive : $\Var(a^tX)=a^t\Sigma_Xa\ge0$. Ainsi $a^tX\sim\Nc(a^tm,\ a^t\Sigma_Xa)$.
\end{remarque}

\begin{theoreme}[fonction caractéristique]\label{C5:fc}
Soit $X$ un vecteur aléatoire de $\R^d$ de carré intégrable, de moyenne $m$ et de covariance $\Sigma_X$. $X$ est un vecteur gaussien \emph{ssi}
\[\varphi_X(u)=\E\bigl(e^{iu^tX}\bigr)=e^{iu^tm-\frac12u^t\Sigma_Xu},\qquad\forall u\in\R^d.\]
Une loi gaussienne est donc entièrement déterminée par $(m,\Sigma_X)$.
\end{theoreme}
\begin{idee}
($\Rightarrow$) $u^tX\sim\Nc(u^tm,u^t\Sigma_Xu)$ donc $\varphi_X(u)=\varphi_{u^tX}(1)$.
($\Leftarrow$) $\varphi_{u^tX}(v)=\varphi_X(vu)=e^{ivu^tm-\frac{v^2}2u^t\Sigma_Xu}$ : c'est la f.c.\ de $\Nc(u^tm,u^t\Sigma_Xu)$.
\end{idee}

\begin{theoreme}[existence et unicité de $\Nc_d(m,K)$]\label{C5:existence}
Soit $m\in\R^d$ et $K$ une matrice $d\times d$ \textbf{symétrique positive} ($x^tKx\ge0$ pour tout $x$).
Il existe une unique loi gaussienne de moyenne $m$ et de covariance $K$, notée $\Nc_d(m,K)$. Les trois conditions
(taille $d\times d$, symétrique, positive) sont nécessaires et suffisantes ; alors les valeurs propres de $K$ sont $\ge0$ et $\det K\ge0$.
\end{theoreme}
\begin{idee}
Existence : $X=(X_1,\dots,X_d)$ avec $X_i\sim\Nc(0,1)$ i.i.d.\ est $\Nc_d(0,I_d)$ (Lemme~\ref{C5:indep-vect}) ; si $K=AA^t$,
$Y=AX+m$ est gaussien de moyenne $m$ et covariance $AI_dA^t=K$ (Lemme~\ref{C5:affine}). Unicité : Thm~\ref{C5:fc}.
\end{idee}

\begin{lemme}[racine carrée de $K$ ; « Cholesky decomposition » dans le cours]\label{C5:racine}
$K$ symétrique positive est diagonalisable en base orthonormée : $K=P\diag(\lambda_1,\dots,\lambda_d)P^t$, $P$ orthogonale, $\lambda_i\ge0$.
Avec $A=P\diag(\sqrt{\lambda_1},\dots,\sqrt{\lambda_d})P^t$, $A$ est symétrique positive et $K=AA^t$.
\end{lemme}
\begin{remarque}
\emph{Hors cours :} la décomposition de Cholesky au sens strict fournit $A$ \emph{triangulaire inférieure} avec $K=AA^t$ ;
pour simuler $\Nc_d(m,K)$, n'importe quel $A$ tel que $K=AA^t$ convient.
\end{remarque}

\begin{horscours}[densité]\label{C5:densite}
Si $K$ est \textbf{inversible}, $\Nc_d(m,K)$ a pour densité sur $\R^d$
\[f(x)=\frac{1}{(2\pi)^{d/2}\sqrt{\det K}}\exp\Bigl(-\frac12(x-m)^tK^{-1}(x-m)\Bigr),\qquad x\in\R^d.\]
Si $\det K=0$, le vecteur vit p.s.\ dans un sous-espace affine strict : \emph{pas de densité} par rapport à $\lambda_d$.
\end{horscours}

\section*{3. Indépendance et covariance}

\begin{theoreme}[blocs]\label{C5:blocs}
Soit $(X_1,\dots,X_m,Y_1,\dots,Y_n)$ un \textbf{vecteur gaussien}. Sont équivalents :
\begin{enumerate}
  \item $X=(X_1,\dots,X_m)$ et $Y=(Y_1,\dots,Y_n)$ sont indépendants ;
  \item $\Cov(X_i,Y_j)=0$ pour tous $1\le i\le m$, $1\le j\le n$ ;
  \item $\Sigma_{(X,Y)}=\begin{pmatrix}\Sigma_X&0\\0&\Sigma_Y\end{pmatrix}$ (diagonale par blocs).
\end{enumerate}
Le résultat s'étend à $k$ blocs. \textbf{Le grand vecteur $(X,Y)$ doit être gaussien}, pas seulement $X$ et $Y$ séparément.
\end{theoreme}
\begin{idee}
$1\Rightarrow2$ toujours. $3\Rightarrow1$ : avec $u=(v,w)$, $u^t\Sigma_{(X,Y)}u=v^t\Sigma_Xv+w^t\Sigma_Yw$, donc
$\varphi_{(X,Y)}(v,w)=\varphi_X(v)\varphi_Y(w)$, ce qui caractérise l'indépendance (loi jointe $=$ produit des lois).
\end{idee}

\begin{corollaire}\label{C5:diag}
Si $X=(X_1,\dots,X_d)$ est un vecteur gaussien, les $X_i$ sont indépendantes $\iff\Sigma_X$ est diagonale.
Et : $X_1,\dots,X_d$ indépendantes de lois $\Nc(m_i,\sigma_i^2)$ $\iff$ $X$ est un vecteur gaussien de moyenne
$(m_1,\dots,m_d)$ et de covariance $\diag(\sigma_1^2,\dots,\sigma_d^2)$.
\end{corollaire}

\begin{liens}
\[\text{gaussiennes indépendantes}\ \Longrightarrow\ \text{vecteur gaussien}\ \Longrightarrow\ \text{composantes gaussiennes}\]
\begin{itemize}
  \item Réciproque 2 fausse : composantes gaussiennes $\not\Rightarrow$ vecteur gaussien. $X\sim\Nc(0,1)$, $a>0$,
  $Y=X\1_{|X|\le a}-X\1_{|X|>a}$ : par symétrie de $\Nc(0,1)$, $Y\sim\Nc(0,1)$ ; mais $X+Y=2X\1_{|X|\le a}$ est bornée,
  non constante, donc non gaussienne : $(X,Y)$ n'est pas un vecteur gaussien.
  \item Réciproque 1 fausse : pour $X\sim\Nc(0,1)$, $(X,X)$ est gaussien mais ses composantes ne sont pas indépendantes ($\Cov(X,X)=1\neq0$).
  \item Dans un vecteur gaussien : composantes non corrélées deux à deux $\iff$ mutuellement indépendantes. Hors vecteur gaussien, seul le sens «~$\Leftarrow$~» subsiste (C3).
\end{itemize}
\end{liens}

\begin{piege}
\begin{itemize}
  \item $\Nc(m,\sigma^2)$ : le second paramètre est la \emph{variance}. $\Var(aX)=a^2\Var X$, pas $a\Var X$.
  \item Pour la somme de gaussiennes, l'\emph{indépendance} (ou le caractère gaussien du vecteur) est indispensable ;
  sinon la somme peut ne pas être gaussienne (contre-exemple ci-dessus).
  \item « Covariance nulle $\Rightarrow$ indépendance » exige que le \emph{vecteur joint} soit gaussien (Thm~\ref{C5:blocs}).
  \item Pour montrer qu'un vecteur est gaussien, l'écrire $AX+b$ avec $X$ gaussien (souvent $X$ à composantes i.i.d.\ $\Nc(0,1)$).
  \item $\Var(X+Y)=\Var X+\Var Y+2\Cov(X,Y)$ : ne pas oublier le terme croisé.
  \item Une loi dégénérée ($\delta_m$, ou $\Sigma$ non inversible) reste gaussienne, mais n'a pas de densité par rapport à $\lambda_d$.
\end{itemize}
\end{piege}

\section*{Exemples types}

\begin{exemple}[somme et différence]\label{C5:ex1}
$X_1,X_2$ i.i.d.\ $\Nc(0,1)$, $U=X_1+X_2$, $V=X_1-X_2$. $(U,V)=A(X_1,X_2)^t$ avec $A=\begin{pmatrix}1&1\\1&-1\end{pmatrix}$ :
vecteur gaussien centré de covariance $AI_2A^t=\begin{pmatrix}2&0\\0&2\end{pmatrix}$. Donc $U\indep V$ (Thm~\ref{C5:blocs}) et
$U,V\sim\Nc(0,2)$.
\end{exemple}

\begin{exemple}[loi d'une combinaison linéaire]\label{C5:ex2}
$X\sim\Nc_2\Bigl(\begin{pmatrix}1\\0\end{pmatrix},\begin{pmatrix}2&1\\1&1\end{pmatrix}\Bigr)$ et $a=(1,-2)^t$.
$a^tm=1$, $a^t\Sigma a=a_1^2\Sigma_{11}+2a_1a_2\Sigma_{12}+a_2^2\Sigma_{22}=2-4+4=2$. Donc $X_1-2X_2\sim\Nc(1,2)$.
\end{exemple}

\end{document}
